\section{Potenzial der LLM-gestützten Aufbereitung für die externe CTI-Nutzung}

Die aktuelle Konzeption der CTI-Reports ist primär auf den Datenaustausch im CSV-Format ausgelegt, welcher für die interne automatisierte Verarbeitung innerhalb der Netzknoten optimiert ist. Ein erhebliches Potenzial bietet jedoch die Aufbereitung dieser Daten mittels Large Language Models (LLMs).

Durch den Einsatz von LLMs könnten die rein maschinenlesbaren CTI-Reports in präzise, natürlichsprachliche Lageberichte überführt werden. Dies ist insbesondere für SOC-Operatoren (Security Operations Center) von hoher Relevanz, da komplexe Bedrohungsszenarien so schneller interpretiert und priorisiert werden können. 